% sections/notations.tex
% --- NOTACIONES Y ACRÓNIMOS (estilo Del Valle) ---
% Del Valle estructura esta sección en dos bloques:
%   1) "Notaciones": prosa que fija las convenciones de símbolos y unidades.
%   2) "Lista de acrónimos": listado a dos columnas (sigla -- significado).
% EDITA el contenido para que refleje las convenciones reales de tu tesis.

\chapter*{Notaciones y acrónimos}
\addcontentsline{toc}{chapter}{Notaciones y acrónimos}
\markboth{Notaciones y acrónimos}{Notaciones y acrónimos}

\section*{Notaciones}

A lo largo del texto se adopta el sistema de unidades $\hbar = 1$, de modo que
energía y frecuencia se expresan en las mismas unidades. % <-- AJUSTA si no aplica a tu caso.
Los operadores se denotan con un acento circunflejo ($\hat{a}$, $\hat{\sigma}$),
y su valor esperado se escribe $\langle \hat{a} \rangle$. Se emplea la notación
habitual de la teoría de conjuntos: $\mathbb{R}$ para los números reales y
$\mathbb{C}$ para los complejos. Los vectores y matrices se escriben en negrita
($\mathbf{x}$, $\mathbf{W}$), y la transpuesta se indica con el superíndice
$^{\mathsf{T}}$.
% AÑADE aquí cualquier otra convención propia de tu trabajo
% (parámetros del láser, hiperparámetros del modelo de aprendizaje, etc.).

\section*{Lista de acrónimos}

\begin{multicols}{2}
\noindent
\begin{description}[leftmargin=5.5em,labelwidth=5em,itemsep=2pt,parsep=0pt]
  \item[\textbf{ML}]   Aprendizaje automático (\textit{Machine Learning})
  \item[\textbf{RL}]   Aprendizaje por refuerzo (\textit{Reinforcement Learning})
  \item[\textbf{NN}]   Red neuronal (\textit{Neural Network})
  \item[\textbf{CNN}]  Red neuronal convolucional
  \item[\textbf{PSO}]  Optimización por enjambre de partículas
  \item[\textbf{GA}]   Algoritmo genético
  \item[\textbf{EA}]   Algoritmo evolutivo
  \item[\textbf{MSE}]  Error cuadrático medio
  \item[\textbf{CW}]   Onda continua (\textit{Continuous Wave})
  \item[\textbf{KLM}]  Anclaje de modos por lente Kerr
  \item[\textbf{SAM}]  Espejo absorbente saturable
  \item[\textbf{SESAM}] Espejo absorbente saturable semiconductor
  \item[\textbf{GVD}]  Dispersión de velocidad de grupo
  \item[\textbf{SPM}]  Automodulación de fase
  \item[\textbf{FWHM}] Anchura a media altura
  \item[\textbf{ODE}]  Ecuación diferencial ordinaria
\end{description}
\end{multicols}
